\documentclass[11pt,a4paper]{article}
\usepackage[margin=2.5cm]{geometry}
\usepackage{graphicx}
\usepackage{booktabs}
\usepackage{hyperref}
\usepackage{amsmath}
\usepackage[utf8]{inputenc}
\usepackage[T1]{fontenc}

\title{Individual Report --- Reinforcement Learning Project\\
\large DQN on \texttt{highway-v0}: Implementation, Infrastructure, and Analysis}
\author{Maxence Rossignol\\
Group 17 --- CentraleSup\'elec, Mention IA}
\date{April 2026}

\begin{document}
\maketitle

\section{Personal Contribution}

My main contributions to the project focused on the \textbf{training infrastructure}, \textbf{experimental pipeline}, and \textbf{cluster deployment}. Specifically:

\begin{itemize}
    \item \textbf{Fair comparison framework}: I converted the DQN training loop from an episode-based to a \textbf{timestep-based} interface, so that both our DQN and the Stable-Baselines3 (SB3) DQN are trained for the same number of environment steps. This ensures the training curves share the same x-axis and the comparison is transparent.
    \item \textbf{Incremental result saving}: I implemented CSV-based logging that saves training rewards and evaluation results after each seed. This allows partial runs to still produce usable data---a design choice that proved critical when our first cluster job was killed mid-training.
    \item \textbf{DCE cluster deployment}: I set up the full training pipeline on CentraleSup\'elec's distributed computing cluster (SLURM-based), including conda environment creation with CUDA 12.4 support, batch job scripts with configurable parameters, and monitoring.
    \item \textbf{Debugging and fixes}: I identified and fixed several issues in the pipeline, including a missing \texttt{Monitor} wrapper for SB3 reward logging, an observation space mismatch caused by \texttt{configure()} requiring a \texttt{reset()} before \texttt{DummyVecEnv} construction, and checkpoint saving during training rather than only at the end.
\end{itemize}

\section{Experiments}

All experiments use the instructor-provided configuration: \texttt{highway-v0} with Kinematics observations ($10 \times 5$ matrix), \texttt{DiscreteMetaAction} (5 actions), 45 vehicles, 4 lanes, 30-second episodes, and normalized rewards.

Both agents share the same core hyperparameters for a fair comparison:
\begin{itemize}
    \item Learning rate: $10^{-3}$, discount factor $\gamma = 0.99$, batch size 64
    \item Replay buffer: 50{,}000 transitions, target network update every 500 steps
    \item Epsilon: linear decay from 1.0 to 0.05 over 50\% of training (DQN) / 30\% (SB3)
    \item Network: MLP with layers $50 \to 256 \to 128 \to 5$ (our DQN); SB3 uses its default \texttt{MlpPolicy}
\end{itemize}

I ran experiments on the DCE cluster (NVIDIA GTX 1080 Ti). A first attempt with 3 seeds was killed by SLURM after exceeding the 5-hour time limit, having completed only seed 42 and 10\% of seed 43's DQN training. The incremental CSV saving ensured no data was lost for the completed seed. This led us to switch to the \texttt{gpu\_prod\_long} partition (48h limit), where the full run---40{,}000 timesteps $\times$ 3 seeds (42, 43, 44) $\times$ 2 agents---completed successfully.

\section{Results and Analysis}

\subsection{Training Curves}

Figure~\ref{fig:training} shows the training curves averaged across 3 seeds, with the shaded area representing $\pm 1$ standard deviation.

\begin{figure}[h]
    \centering
    \includegraphics[width=0.85\textwidth]{../results_distant_v2/figures/training_curves.png}
    \caption{Training curves (mean $\pm$ std across seeds 42, 43, 44), 40{,}000 timesteps.}
    \label{fig:training}
\end{figure}

Key observations:
\begin{itemize}
    \item Our DQN reaches high reward levels ($\sim$17--20) within the first 5{,}000 timesteps, while SB3 takes roughly 20{,}000--25{,}000 timesteps to reach comparable performance. This is likely because SB3 delays learning for the first 1{,}000 steps (\texttt{learning\_starts} parameter) and uses a more conservative exploration schedule.
    \item Our DQN shows tight variance across seeds (narrow shaded region), indicating \textbf{robust and reproducible} learning. SB3's variance is much wider, especially between 10{,}000 and 30{,}000 timesteps, reflecting inconsistent learning across seeds.
    \item Both agents converge toward similar performance ($\sim$20) by the end of training, but SB3 remains noisier.
\end{itemize}

\subsection{Evaluation}

Table~\ref{tab:eval} summarizes the evaluation results over 50 episodes per seed with deterministic policies (no exploration).

\begin{table}[h]
    \centering
    \caption{Evaluation results (50 episodes per seed, deterministic policy).}
    \label{tab:eval}
    \begin{tabular}{lccccc}
        \toprule
        Agent & Seed & Mean Reward & Std & Crashes & Crash Rate \\
        \midrule
        DQN (ours) & 42 & 20.68 & 2.70 & 6/50 & 12\% \\
        DQN (ours) & 43 & 20.33 & 2.60 & 2/50 & 4\% \\
        DQN (ours) & 44 & 20.53 & 2.29 & 4/50 & 8\% \\
        \textbf{DQN (ours)} & \textbf{all} & \textbf{20.51} & \textbf{2.54} & \textbf{12/150} & \textbf{8\%} \\
        \midrule
        SB3 DQN & 42 & 19.52 & 3.47 & 13/50 & 26\% \\
        SB3 DQN & 43 & 16.11 & 6.00 & 22/50 & 44\% \\
        SB3 DQN & 44 & 10.77 & 7.35 & 42/50 & 84\% \\
        \textbf{SB3 DQN} & \textbf{all} & \textbf{15.47} & \textbf{6.85} & \textbf{77/150} & \textbf{51\%} \\
        \bottomrule
    \end{tabular}
\end{table}

The results are striking. Our DQN significantly outperforms SB3 on all metrics: higher mean reward (20.51 vs 15.47), lower variance (2.54 vs 6.85), and dramatically fewer crashes (8\% vs 51\%). Notably, SB3's performance is highly seed-dependent---seed 42 performs reasonably (26\% crashes) while seed 44 crashes in 84\% of episodes, suggesting that SB3's learned policy is fragile and depends heavily on the exploration trajectory during training.

\subsection{Failure Analysis}

Crash episodes in our DQN follow two patterns:
\begin{itemize}
    \item \textbf{Early crashes} (steps 3--9, e.g., reward 6.56 at step 8): the agent attempts a lane change immediately after reset into an occupied lane. These represent rare but catastrophic misjudgments of the initial traffic configuration.
    \item \textbf{Late crashes} (steps 21--29, e.g., reward 19.51 at step 25): the agent survives most of the episode but collides near the end, typically during a lane change into fast-moving traffic. The crash speed ($\sim$13--19 m/s) suggests these are mid-maneuver collisions where a faster vehicle in the target lane caught up during the lane change.
\end{itemize}

SB3's failure modes are more diverse and severe. Many crashes occur mid-episode at high speeds (up to 28 m/s), and for seed 44, the agent appears to have learned an unstable policy that fails in the majority of traffic configurations.

\section{What Worked and What Failed}

\textbf{What worked:}
\begin{itemize}
    \item The timestep-based unification made the comparison genuinely fair and the training curves directly readable on the same x-axis.
    \item Incremental CSV saving paid off immediately when our first cluster job timed out. Without it, 2+ hours of compute would have been lost entirely.
    \item Our from-scratch DQN significantly outperforms SB3's optimized implementation on this benchmark, which validates both the implementation and the hyperparameter choices.
\end{itemize}

\textbf{What failed:}
\begin{itemize}
    \item \textbf{SB3 reward logging}: the \texttt{RewardLoggerCallback} silently produced empty lists because the environment was not wrapped in a \texttt{Monitor}. This went undetected until we noticed the SB3 training curve was missing from our first plots. The fix was a single line (\texttt{Monitor(env)}).
    \item \textbf{Compute budget estimation}: I underestimated how long 40{,}000 timesteps $\times$ 3 seeds would take on the cluster ($\sim$14 hours), leading to the first job being killed at the SLURM time limit. This was solved by switching to the \texttt{gpu\_prod\_long} partition (48h limit).
    \item \textbf{Observation space bug}: SB3's \texttt{DummyVecEnv} reads the observation space at construction time, before \texttt{configure()} takes effect. The fix was to call \texttt{env.reset()} after \texttt{configure()} so the observation space updates from $(5, 5)$ to $(10, 5)$.
\end{itemize}

\section{Limitations}

\begin{itemize}
    \item \textbf{No hyperparameter tuning} was performed. The learning rate, network architecture, epsilon schedule, and buffer size were set to reasonable defaults without systematic optimization. SB3's poor performance may partly be due to its \texttt{learning\_starts=1000} and \texttt{exploration\_fraction=0.3} settings, which could be suboptimal for this environment.
    \item The \textbf{extension task} has not yet been completed. Future work could explore Double-DQN and Prioritized Experience Replay as natural improvements to the current DQN implementation.
    \item The failure analysis is qualitative. A more rigorous approach would involve recording the full observation trajectory leading to each crash to identify systematic failure patterns.
    \item While 3 seeds is the minimum required, more seeds (e.g., 5--10) would provide tighter confidence intervals, especially for SB3 which shows high seed-to-seed variance.
\end{itemize}

\end{document}
